% Chapter 02.1: কমিউনিকেশন সিস্টেম ও ব্যান্ডউইডথ
\section{কমিউনিকেশন সিস্টেম ও ব্যান্ডউইডথ}

\subsection{কমিউনিকেশন সিস্টেমের ধারণা}

\begin{itemize}
\item কমিউনিকেশন কী
\item কমিউনিকেশন সিস্টেম কী
\item তথ্য প্রেরণ ও গ্রহণের ধারণা
\item ডেটা ও সিগন্যালের সম্পর্ক
\item Analog signal ও Digital signal
\end{itemize}

\subsection{কমিউনিকেশন মডেল}

\begin{itemize}
\item Source System
\item Transmitter
\item Transmission System / Medium
\item Receiver
\item Destination System
\end{itemize}

\subsection{ডেটা কমিউনিকেশনের ধারণা}

\begin{itemize}
\item Data communication কী
\item ডেটা প্রেরণ ও গ্রহণ
\item ডেটা কমিউনিকেশনের প্রয়োজনীয়তা
\item Sender / Source, Receiver / Destination, Medium, Message, Protocol
\end{itemize}

\subsection{ব্যান্ডউইডথ}

\begin{itemize}
\item ব্যান্ডউইডথের ধারণা
\item ডেটা প্রেরণের ক্ষমতা
\item bps, Kbps, Mbps, Gbps
\item Narrow Band, Voice Band, Broadband
\end{itemize}
